
\subsubsection{3.1 Hypothesis: Preference Entropy Collapse Under RLHF}

We hypothesize that extended RLHF training induces preference entropy collapse faster than task performance improves. Specifically:

\textbf{Core Statement:} Under RLHF training on preference datasets, if models undergo extended alignment optimization (10K-20K training steps), then preference distribution entropy collapses faster than task performance improves (inflection point where dH/dP accelerates), because users habituate to model output style and converge on preferences even for subjective tasks where diversity is legitimate.

\textbf{Causal Mechanism (4 steps):}
\begin{enumerate}
\item Base models $\rightarrow$ high-variance outputs $\rightarrow$ diverse preferences $\rightarrow$ high entropy (baseline)
\item Early RLHF (1K-5K steps) $\rightarrow$ quality improvement $\rightarrow$ entropy decreases proportionally to performance (justified reduction)
\item Extended RLHF (10K-20K steps) $\rightarrow$ performance saturates $\rightarrow$ entropy reduction accelerates beyond performance gains (inflection point)
\item Task stratification $\rightarrow$ subjective tasks show greater entropy collapse than objective tasks (monoculture evidence)
\end{enumerate}

\textbf{Key Assumptions:}
\begin{itemize}
\item \textbf{A1:} Preference entropy proxies critical evaluation capacity (higher H = diverse judgments)
\item \textbf{A5:} Existing RLHF datasets contain raw preference distributions (not just aggregated win rates)
\end{itemize}

This paper validates Assumption A5 via sub-hypothesis H-E1. The full causal mechanism (steps 2-4) requires this prerequisite to succeed.

\subsubsection{3.2 Sub-Hypothesis H-E1: Preference Entropy Measurement Feasibility}

\textbf{Statement:} Under RLHF preference datasets (Anthropic-HH), if we have access to raw pairwise comparison data, then we can compute Shannon entropy H = -$\Sigma$ p\_i log(p\_i) for preference distributions, because the dataset publishes response frequency counts rather than only aggregated win rates.

\textbf{Verification Protocol:}
\begin{enumerate}
\item Download Anthropic-HH dataset (160K+ pairwise comparisons)
\item Sample n=100 prompts with $\geq 5$ comparisons each (seed=1)
\item Group examples by prompt, aggregate preference distributions
\item Compute Shannon entropy H for each prompt's preference distribution
\item Validate entropy range [0, ln(2)] for binary comparisons
\item Report success rate (\% of prompts with computable entropy) and variance (std of entropy values)
\end{enumerate}

\textbf{Success Criteria (MUST\_WORK gate):}
\begin{itemize}
\item \textbf{Primary:} Entropy computable for $\geq 95$\% of sampled prompts
\item \textbf{Secondary:} Entropy variance > 0 (not all constant values)
\end{itemize}

\textbf{Rationale:} If entropy cannot be computed from existing datasets (A5 violated), the hypothesis mechanism becomes untestable without new data collection. H-E1 validates the foundational measurement assumption.

\subsubsection{3.3 Entropy Measurement}

\textbf{Shannon Entropy Definition:}

$H = -\Sigma p_i log(p_i)$

where p\_i = proportion of preferences for option i, n = number of response options per prompt, natural log (base e) for units in nats.

\textbf{For Binary Pairwise Comparisons:}
\begin{itemize}
\item n = 2 (chosen vs rejected)
\item Theoretical range: [0, ln(2)] $\approx$ [0, 0.693] nats
\item Maximum entropy (uniform): H = ln(2) when p\_chosen = p\_rejected = 0.5
\item Minimum entropy (consensus): H = 0 when all prefer one option
\end{itemize}

\textbf{Expected Behavior:}
\begin{itemize}
\item High diversity prompts: H > 0.4 nats
\item Low diversity prompts: H < 0.3 nats
\item Variance across prompts: std(H) > 0.1 nats indicates meaningful spread
\end{itemize}

\subsubsection{3.4 Dataset: Anthropic-HH}

\textbf{Source:} Anthropic Helpful \& Harmless (HH) dataset \citep{bai2022constitutional}
\begin{itemize}
\item \textbf{Repository:} https://github.com/anthropics/hh-rlhf
\item \textbf{Size:} 160,800 training examples
\item \textbf{Format:} Each example contains \texttt{chosen} (preferred response) and \texttt{rejected} (less-preferred response) full conversation texts
\item \textbf{Collection:} Single-annotator pairwise comparisons
\end{itemize}

\textbf{Sampling Strategy:}
\begin{itemize}
\item Sample n=100 prompts from training split (seed=1)
\item Filter: require $\geq 5$ examples per prompt group
\item Grouping: extract prompt from first 200 characters of chosen response text
\end{itemize}

\textbf{Justification:}
\begin{enumerate}
\item Public availability (reproducible validation)
\item Large scale (160K+ comparisons)
\item RLHF relevance (used in Constitutional AI)
\item Raw data access (individual comparison records)
\end{enumerate}

\subsubsection{3.5 Dataset Structure Documentation}

We document two dataset format categories relevant for entropy measurement:

\begin{table*}[htbp]
\centering
\resizebox{\dimexpr 2\columnwidth+\columnsep\relax}{!}{%
\begin{tabular}{lllll}
\toprule
Format Type & Structure & Annotation Cost & Entropy Measurable & Example Dataset \\
\midrule
**Pairwise Unique Responses** & Each comparison evaluates different response candidates (A1 vs B1, A2 vs B2) & Low (1 annotator per comparison) & **No** --- cannot aggregate into distributions over fixed options & Anthropic-HH (validated), WebGPT (inferred), InstructGPT (inferred) \\
**Multi-Annotator Fixed Responses** & Multiple annotators rate same response candidates per prompt (N voters on {R1, R2, ...}) & High (N annotators $\times$ M responses) & **Yes (inferred)** --- distribution over fixed response set & OpenAI Summarization (inferred), Chatbot Arena (inferred) \\
\bottomrule
\end{tabular}
}
\end{table*}

\textbf{Note:} Multi-annotator measurability is theory-based (positive case untested empirically). Table documents cost-measurement trade-off observed in Anthropic-HH with format-based inference for other datasets.

\textbf{Key Insight:} Pairwise-unique format optimizes annotation cost by assigning different response pairs to each comparison. This prevents per-prompt entropy aggregation because:

\begin{itemize}
\item \textbf{Intended aggregation:} ''For prompt P, what \% prefer response A vs B vs C?'' (requires fixed {A, B, C})
\item \textbf{Actual structure:} ''Comparison 1: chosen=A1, rejected=B1; Comparison 2: chosen=A2, rejected=B2'' (no shared response set)
\item \textbf{Tautological result:} Treating each comparison as ''1 vote for chosen category'' yields 50/50 split $\rightarrow$ H = ln(2) for all prompts
\end{itemize}

Multi-annotator-fixed format avoids this by fixing response candidates per prompt, enabling variance across prompts.

\subsubsection{3.6 Implementation}

\textbf{Tools:}
\begin{itemize}
\item Dataset loading: Hugging Face \texttt{datasets} library
\item Entropy computation: \texttt{scipy.stats.entropy}
\item Visualization: matplotlib
\end{itemize}

\textbf{Code Structure:}
\begin{verbatim}
from datasets import load_dataset
from scipy.stats import entropy
import numpy as np

class PreferenceEntropyAnalyzer:
    def compute_entropy_for_prompt(self, prompt_examples):
        chosen_count = len(prompt_examples)
        rejected_count = len(prompt_examples)
        preference_counts = np.array([chosen_count, rejected_count])
        return entropy(preference_counts, base=np.e)
\end{verbatim}

\textbf{Validation Checks:}
\begin{itemize}
\item All entropy values in range [0, ln(2)] nats
\item No NaN or Inf values
\item Deterministic results (seed=1 reproducibility)
\end{itemize}

